%% LyX 2.4.4 created this file.  For more info, see https://www.lyx.org/.
%% Do not edit unless you really know what you are doing.
\documentclass[english]{article}
\usepackage[T1]{fontenc}
\usepackage[latin9]{inputenc}
\usepackage{units}
\usepackage{enumitem}
\usepackage{amsmath}
\usepackage{amsthm}
\usepackage{amssymb}
\usepackage{cancel}
\usepackage{geometry}
\geometry{verbose,tmargin=1in,bmargin=1in,lmargin=1in,rmargin=1in}

\makeatletter
%%%%%%%%%%%%%%%%%%%%%%%%%%%%%% Textclass specific LaTeX commands.
\numberwithin{equation}{section}
\numberwithin{figure}{section}
\newlength{\lyxlabelwidth}      % auxiliary length 

%%%%%%%%%%%%%%%%%%%%%%%%%%%%%% User specified LaTeX commands.
\usepackage{fancyhdr}
\usepackage{lastpage}
\pagestyle{fancy}
\usepackage{enumitem}
\usepackage{ifthen}
\everymath=\expandafter{\the\everymath\displaystyle}
%\DeclareMathSizes{10}{10}{10}{10}
\usepackage{multicol}

\usepackage{tikz}
\usetikzlibrary{patterns}
\usetikzlibrary{plotmarks}
\usepackage{pgfplots}
\pgfplotsset{compat=newest}

\usepackage{advdate}    % Advancing/saving dates
\usepackage{datetime}   % Dates formatting
\usepackage{datenumber} % Counters for dates
\usdate
% Added by lyx2lyx
\setlength{\parskip}{\smallskipamount}
\setlength{\parindent}{0pt}

\makeatother

\usepackage{babel}
\begin{document}
\renewcommand{\author}{Dr.\,James Ashe}

\newcommand{\classNum}{335}

\newcommand{\classSec}{A}

\newcommand{\examType}{Practice Midterm \textbf{due: \formatdate{4}{3}{2014}}}

\renewcommand{\date}{}

\newcommand{\Name}{Solutions} %name:

%%First page's header%%%%%%%%%%%%%%%%%%%%%%
\fancypagestyle{firststyle} %%header settings for first page
{
	\fancyhead[L]{MTH \classNum{}(\classSec{}) \author{}\\
	\textbf{\examType{}} \date{}}
	\fancyhead[C]{\textbf{\Name}}
	\setlength{\headsep}{14pt}
}
%%%%%%%%%%%%%%%%%%%%%%%%%%%%%%%%%%%%%%%%%%

%%Next page's header%%%%%%%%%%%%%%%%%%%%%%%
\setlist{nolistsep}
\lhead{\classNum{}(\classSec{})} 
\chead{\textbf{}} 
\cfoot{\thepage{}~of~\pageref{LastPage}} 
\setlength{\headsep}{5pt} 
%%%%%%%%%%%%%%%%%%%%%%%%%%%%%%%%%%%%%%%%%%%

%%Various settings; see the preamble for other settings%%%%%
%\setenumerate[0]{leftmargin=12pt,itemindent=0pt,labelwidth=0pt}
\setlist{topsep=.5pt}
\renewcommand{\labelenumi}{\textbf{\arabic{enumi}.}}
\renewcommand{\labelenumii}{\textbf{\alph{enumii}.}}
\thispagestyle{firststyle}
%%%%%%%%%%%%%%%%%%%%%%%%%%%%%%%%%%%%%%%%%%

\newcommand{\Dspace}{.5cm}
\begin{enumerate}[resume]
\item Determine the following

\begin{enumerate}
\item $17\bmod2$

(answers may vary) $17\bmod2=1$. More generally, $17\bmod2=2k+1$
for any $k\in\left\{ 0,1,2,3,\dots\right\} $ since $2|17-(2k+1)$.
\item $\mbox{lcm}(7^{4000}\cdot5^{2}\cdot2^{4},7^{4}\cdot5^{1000}\cdot2)$

$\mbox{lcm}(7^{4000}\cdot5^{2}\cdot2^{4},7^{4}\cdot5^{1000}\cdot2)=7^{4000}\cdot5^{1000}\cdot2^{4}$
\item $17=x\cdot2-y\cdot15$. What is $\gcd(x,y)$?

The $\gcd(x,y)$ is the least positive integer representable in the
form $x\cdot s+y\cdot t$. All others are \emph{multiples }of $\gcd(x,y)$,
e.g., $2\cdot1+6\cdot1=8$, $\gcd(2,6)=2$ and $2|8$. \\
Since $x$ and $y$ can be written as a linear combination summing
to $17$, the $\gcd(x,y)=\mbox{either }1\mbox{ \textbf{or} }17$.\vfill
\end{enumerate}
\item Complete the following

\begin{enumerate}
\item Write the given set in tabular form: $A=\{x\mid x\in\mathbb{N}\mbox{ and }x\leq5\}$

$A=\left\{ 1,2,3,4,5\right\} $\vfill
\item Write the given set in set-builder form: $A=\left\{ \dots-2,0,2,4\right\} $

$A=\left\{ x|x\mbox{ is an even integer and }x\leq4\right\} $ or
$A=\left\{ x|x=2n\leq4\mbox{ and \ensuremath{n\in\mathbb{Z}}}\right\} $\vfill
\end{enumerate}
\item Let $A=\left\{ 1,2,3\right\} $, $B=\left\{ 2,4,6\right\} $, $C=\left\{ 4,5,6\right\} $.

\begin{enumerate}
\item Find $A\cup C$

$A\cup C=\left\{ 1,2,3,4,5,6\right\} $
\item Find $A\cap B$

$A\cap B=\left\{ 2\right\} $
\item Find $A\cap C$

$A\cap C=\emptyset$
\item Find $\left(A\cup C\right)\cap B$

$\left(A\cup C\right)\cap B=\left\{ 1,2,3,4,5,6\right\} \cap B=\left\{ 2,4,6\right\} $\vfill
\end{enumerate}
\item Let $A$ and $B$ be two sets. Prove that $A\cap B\subseteq A$. 

\begin{proof}
Let $x\in A\cap B$. Then $x\in B$ and $x\in A$. Thus $x\in A$
and $A\cap B\subseteq A$.\vfill
\end{proof}
\item Let $A\neq\emptyset$ be a subset of $B$. Prove that $A\cup B\neq A$.

\begin{proof}
Let $A=\left\{ a\right\} $ and $B=\left\{ a,b\right\} $. Then $A\subseteq B$,
$A\neq\emptyset$, but $\left\{ a,b\right\} =A\cup B\neq A=\left\{ a\right\} $.
\vfill
\end{proof}
\item Give the definition of the following:

\begin{enumerate}
\item A \emph{function}

A \emph{function} is a relation that pairs element from a set $A$
to exactly one element in a set $B$.
\item An \emph{onto function}.

A function $f:A\rightarrow B$ is \emph{onto }if for every $b\in B$
there exists an $a\in A$ such that $f(a)=b$.\vfill
\end{enumerate}
\item Consider the map $f:\mathbb{R}\times\mathbb{R}\rightarrow\mathbb{R}$
defined as $f(x,y)=xy$. Then $f$ is:

\begin{enumerate}
\item [(a)] an onto function.
\item [(b)] a 1-1 function.
\item [(c)] a bijection.
\item [(d)] none of the above.

$f$ is onto since for every $b\in\mathbb{R}$, $f(b,1)=b\cdot1$.
$f$ is not $1-1$ since $f(-2,-2)=4$ and $f(2,2)=4$. \\
So the answer is \textbf{(a)}.\vfill
\end{enumerate}
\item Let $f:A\rightarrow B$ be a function, $A=\left\{ 1,2,3\right\} $,
and $B=\left\{ 1,2\right\} $. 

\begin{enumerate}
\item Can $f$ be 1-1? Explain.

No. $\left|A\right|>\left|B\right|$ so $f$ cannot be $1-1$ by the
Pigeon Hole Theorem. 
\item Show that $f$ can be onto.

(answers can vary) Let $f(1)=1$, $f(2)=2$, and $f(3)=1$. \\
Note that $f$ does not have to be onto, e.g, let $f$ be the constant
function, $f(x)=1$.\vfill
\end{enumerate}
\item Define the following relation on the real number set: $x\mathcal{R}y\iff xy>0$.
Show the following:

\begin{enumerate}
\item $\mathcal{R}$ is \emph{not }reflexive.

\begin{proof}
Let $x=0$. Then $x\cdot x=0\not>0$. So $\mathcal{R}$ is not reflexive.
\end{proof}
\item $\mathcal{R}$ is transitive.

\begin{proof}
Let $x,y,z\in\mathbb{R}$. If $x\cdot y>0$ and $y\cdot z>0$ then
$\nicefrac{xy}{y}=x>0$ and $\nicefrac{yz}{y}=z>0$. Thus $x\cdot z>0$.
So $\mathcal{R}$ is transitive.
\end{proof}
\item $\mathcal{R}$ is \emph{not }an equivalence relation.

\begin{proof}
$\mathcal{R}$ is not reflexive by part (1). So $\mathcal{R}$ is
not an equivalence relation.\vfill
\end{proof}
\end{enumerate}
\item Give an example of a relation on a set which is an equivalence relation.

(answers may vary) Let $x\mathcal{R}y\iff x=y$. \\
$x=x$ so $\mathcal{R}$ is reflexive. $x=y\Rightarrow y=x$ so $\mathcal{R}$
is symmetric. If $x=y$ and $y=z$ then $x=y=z$ so $\mathcal{R}$
is transitive. Thus $\mathcal{R}$ is an equivalence relation. \vfill
\end{enumerate}

\end{document}
